\subsection{Simplex4TwIS}

\begin{frame}
  \frametitle{\currentsectionname}

  \begin{figure}
    \tikzstyle{s4dArrow}=[->, line width=1.5mm, s4d-blue]
\centering
\resizebox{.95\textwidth}{!}{%
  \begin{tikzpicture}
    \node (s4dImg) at (8, -6)
    [inner sep=0, outer sep=0]
    {\includegraphics[height=4cm]{assets/s4t.png}};

    \node (s4d) at ($(s4dImg.center) + (0.15,0)$)
    [align=center, minimum width=3.1cm, minimum height=4.25cm]
    {};

    \node (s4dReality) at (s4d.north east)
    [anchor=south east, align=center]
    {\huge Realitäts-\\\huge modell};

    \node (s4dImport) at ($(s4d.west) - (5,0)$)
    [anchor=east, align=center]
    {\huge Import};

    \node (s4dScenarios) at ($(s4d.east) + (5,0)$)
    [anchor=west, align=center, minimum width=3.5cm]
    {\huge Simplex\\\huge Szenarios};

    \node (s4dService) at (s4dScenarios |- s4dReality.base)
    [anchor=base, align=center, minimum width=3.5cm]
    {\huge Simplex\\\huge Service};

    \draw [s4dArrow]
    (s4dImport) edge node (importArrow) {} (s4d)
    (s4d) edge node (scenarioArrow) {} (s4dScenarios)
    (s4dReality) edge (s4dService)  (s4dScenarios) edge (s4dService);

    \node (block1) at ($(importArrow.center) - (0,2)$)
    [align=center]
    {\Large Konvertierung zu\\\Large Realitätsmodell};

    \node (block2) at ($(scenarioArrow.center) - (0,2)$)
    [align=center]
    {\Large Zusammenstellung von\\\Large Anwendungsfällen};

    \node (importCircle) at (s4dImport.center)
    [shape=circle, minimum width=7cm]
    {};

    \node (i1) at (importCircle.140) {\Large{}.shp};
    \node (i2) at (importCircle.160) {\Large{}.json};
    \node (i3) at (importCircle.180) {\Large{}.csv};
    \node (i4) at (importCircle.200) {\Large{}API};
    \node (i5) at (importCircle.220) {\Large{}\dots};

    \draw [s4dArrow, line width=.5mm]
    (i1) edge (s4dImport)
    (i2) edge (s4dImport)
    (i3) edge (s4dImport)
    (i4) edge (s4dImport)
    (i5) edge (s4dImport);

  \end{tikzpicture}
}

    \caption{Datenverarbeitung in Simplex4TwIS.}
  \end{figure}

  \Note{
      \item betrachten Simplex4TwIS aus Sicht der Daten, sieht so aus
      \item Aus verschiedenen Dateiformaten
  }

\end{frame}

\begin{frame}

  \begin{itemize}
    \item \textbf{Realitätsmodell}: Objektorientiertes und graphenbasiertes Datenmodell.
      \begin{itemize}
        \item \textbf{Harmonisiert}: Objekte, Attribute und Verbindungen besitzen immer gleiche Datenstrukturen.
        \item \textbf{Atomar}: Jede Dateneinheit wird als Objektklasse abgebildet.
      \end{itemize}
  \end{itemize}

  \Note{
    \item Standardfelder: Schlüssel immer gleich, können jedoch durch zusätzliche Metadaten wie Name und Kommentar weiter beschrieben werden
    \item zusätzliche Informationen über Sachattribute, welche in den Metadaten über Standardfelder beschrieben werden
    \item Datenmodell orientiert sich nicht an fachspezifischen Anforderungen
    \item anwendungsspezifische Auswertungen basierend auf Realitätsmodell führen dazu dass keine fachspezifischen Modelle mehr benötigt werden
  }

\end{frame}
